\documentclass[11pt,a4paper]{article}
\usepackage[italian]{babel}
\usepackage[margin=2.5cm]{geometry}
\usepackage{parskip}
\usepackage{amsmath, amssymb, amsthm}
\usepackage{booktabs}
\usepackage{array}
\usepackage{xcolor}
\usepackage{needspace}
\usepackage{enumitem}
\usepackage{listings}
\usepackage{tikz}
\usepackage[most]{tcolorbox}
\usepackage[hidelinks]{hyperref}
\usetikzlibrary{decorations.pathreplacing, shapes.geometric, arrows.meta, positioning}

% Niente vedove/orfane: i paragrafi non lasciano righe isolate a fine/inizio pagina
\widowpenalty=10000
\clubpenalty=10000

% Elenchi più compatti
\setlist{itemsep=2pt, parsep=0pt, topsep=4pt, partopsep=0pt}

% --- Riquadri con bordo nero, senza numero (si spezzano tra due pagine) ---
% Uso: \begin{definizione} ... \end{definizione}
%      \begin{teorema}[Nome] ... \end{teorema}  (il nome è facoltativo)
\tcbset{nota/.style={enhanced jigsaw, breakable, colback=white,
  colframe=black, boxrule=0.5pt, sharp corners}}
\NewTColorBox{definizione}{}{nota,
  before upper={\textbf{Definizione.}\ }}
\NewTColorBox{teorema}{o}{nota,
  before upper={\textbf{Teorema\IfValueT{#1}{ (#1)}.}\ }}
\NewTColorBox{proposizione}{o}{nota,
  before upper={\textbf{Proposizione\IfValueT{#1}{ (#1)}.}\ }}
\NewTColorBox{proprieta}{o}{nota,
  before upper={\textbf{Proprietà\IfValueT{#1}{ (#1)}.}\ }}
\NewTColorBox{assioma}{o}{nota, boxrule=1pt,
  before upper={\textbf{Assioma\IfValueT{#1}{ (#1)}.}\ }}

% --- Ambienti semplici (senza riquadro) ---
% Il titolo sta in cima al blocco, anche se il contenuto inizia con un riquadro o
% con codice; e non resta mai da solo in fondo a una pagina.
% Uso: \begin{esempio}[Titolo facoltativo] ... \end{esempio}
\newcommand{\newrunin}[2]{%
  \NewDocumentEnvironment{#1}{o}{%
    \par\Needspace{4\baselineskip}\addvspace{\medskipamount}\noindent
    \textbf{#2}\IfValueT{##1}{ (##1)}\textbf{.}\ \ignorespaces}%
  {\par\addvspace{\medskipamount}}}
\newrunin{esempio}{Esempio}
\newrunin{esercizio}{Esercizio}
\newrunin{osservazione}{Osservazione}

% V e F nelle tavole di verità
\newcommand{\V}{\textsf{V}}
\newcommand{\F}{\textsf{F}}

% --- Codice C ---
% Uso: \begin{codice} ... \end{codice}      codice C numerato
%      \begin{comando} ... \end{comando}    comandi da terminale
%      \begin{risultato} ... \end{risultato} output di un programma
%      \lstinline|printf("ciao");|           codice dentro al testo
\lstdefinestyle{c}{
  language=C,
  basicstyle=\ttfamily\small,
  keywordstyle=\color{blue}\bfseries,
  commentstyle=\color{gray}\itshape,
  stringstyle=\color{red!70!black},
  numbers=left,
  numberstyle=\tiny\color{gray},
  numbersep=8pt,
  columns=flexible,
  keepspaces=true,
  breaklines=true,
  showstringspaces=false,
  tabsize=4,
  morekeywords={include, define},
  % lettere accentate nei commenti
  literate={à}{{\`a}}1 {è}{{\`e}}1 {é}{{\'e}}1 {ì}{{\`i}}1
           {ò}{{\`o}}1 {ù}{{\`u}}1 {È}{{\`E}}1 {À}{{\`A}}1
}
\lstset{style=c}
\newtcblisting{codice}{listing only, nota, left=6mm,
  listing options={style=c}}
\newtcblisting{comando}{listing only, nota,
  listing options={basicstyle=\ttfamily\small, numbers=none, language={},
    columns=flexible, keepspaces=true, breaklines=true}}
\newtcblisting{risultato}{listing only, enhanced jigsaw, breakable,
  colback=black!5, colframe=black, boxrule=0.5pt, sharp corners,
  listing options={basicstyle=\ttfamily\small, numbers=none, language={},
    columns=flexible, keepspaces=true, breaklines=true}}

% --- Diagrammi di flusso (simboli standard) ---
\tikzset{
  terminale/.style = {ellipse, draw, minimum width=2.5cm, minimum height=0.9cm},
  io/.style        = {trapezium, trapezium left angle=70, trapezium right angle=110,
                      draw, minimum width=2.5cm, minimum height=0.9cm},
  processo/.style  = {rectangle, draw, minimum width=2.5cm, minimum height=0.9cm},
  decisione/.style = {diamond, draw, aspect=2, minimum width=2.5cm},
  freccia/.style   = {thick, -{Stealth}}
}

\title{Limiti di esponenziali, logaritmi e trigonometriche}
\author{Marco Cerbella}
\date{Lezione 13 -- 5 novembre 2026}

\begin{document}
\maketitle

\section{Limiti delle funzioni esponenziali e logaritmiche}

Sia $a > 0$, $a \neq 1$. Valgono
\[
\lim_{x\to+\infty} a^x = \begin{cases} +\infty & a > 1 \\ 0 & 0 < a < 1 \end{cases}
\qquad
\lim_{x\to-\infty} a^x = \begin{cases} 0 & a > 1 \\ +\infty & 0 < a < 1 \end{cases}
\]
\[
\lim_{x\to+\infty} \log_a x = \begin{cases} +\infty & a > 1 \\ -\infty & 0 < a < 1 \end{cases}
\qquad
\lim_{x\to 0^+} \log_a x = \begin{cases} -\infty & a > 1 \\ +\infty & 0 < a < 1 \end{cases}
\]
Il grafico di $a^x$ passa per $(0, 1)$, quello di $\log_a x$ per $(1, 0)$.

\begin{esempio}[Giustificazione dei limiti dell'esponenziale]
Se $0 < a < 1$ si pone $a = \dfrac{1}{b}$ con $b > 1$, e quindi $a^x = \left(\dfrac{1}{b}\right)^x = \dfrac{1}{b^x}$.
\begin{itemize}
  \item $0 < a < 1$, $x \to +\infty$: $a^x = \dfrac{1}{b^x} \to 0$, perché $b^x \to +\infty$ ($b > 1$).
  \item $a > 1$, $x \to -\infty$: si pone $x = -z$ con $z > 0$ (e $z \to +\infty$); allora $a^x = a^{-z} = \dfrac{1}{a^z} \to 0$.
  \item $0 < a < 1$, $x \to -\infty$: $a = \dfrac{1}{b}$, $x = -z$, $z > 0$, da cui $a^x = \left(\dfrac{1}{b}\right)^{-z} = \dfrac{1}{b^{-z}} = b^z \to +\infty$.
\end{itemize}
\end{esempio}

\begin{esempio}[Giustificazione dei limiti del logaritmo]
Sia $\log_a x = c \iff a^c = x$.
\begin{itemize}
  \item Se $a > 1$ e $x \to +\infty$, allora $c \to +\infty$.
  \item Se $0 < a < 1$ e $x \to +\infty$: posto $a = \dfrac{1}{b}$ con $b > 1$, si ha $\log_{1/b} x = c \iff \left(\dfrac{1}{b}\right)^c = x$. Il limite $x \to +\infty$ è possibile solo se $c < 0$, quindi $c \to -\infty$: infatti se $c < 0$, $c = -d$ con $d > 0$, e $\left(\dfrac{1}{b}\right)^c = \left(\dfrac{1}{b}\right)^{-d} = b^d$.
  \item Se $a > 1$ e $x \to 0^+$: $c < 0$ e $c \to -\infty$.
  \item Se $0 < a < 1$ e $x \to 0^+$: $c > 0$ e $c \to +\infty$. Infatti $\left(\dfrac{1}{b}\right)^c = x \to 0$, cioè $\dfrac{1}{b^c} \to 0$, cioè $b^c \to +\infty$, cioè $c \to +\infty$.
\end{itemize}
\end{esempio}

\section{Esempi di limiti con somme}

\begin{esempio}
\begin{itemize}
  \item $\displaystyle\lim_{x\to+\infty}\left(\frac{1}{x} + e^x\right) = +\infty$, perché $\frac{1}{x} \to 0$ e $e^x \to +\infty$ ($e > 1$).
  \item $\displaystyle\lim_{x\to-\infty}\left(\frac{1}{x} + e^x\right) = 0$, perché $\frac{1}{x} \to 0$ e $e^x \to 0$.
  \item $\displaystyle\lim_{x\to+\infty}\left(\frac{1}{x} + e^{-x}\right) = 0$, perché $\frac{1}{x} \to 0$ e $e^{-x} \to 0$.
  \item $\displaystyle\lim_{x\to+\infty}\left(x + \ln x\right) = +\infty$, perché $x \to +\infty$ e $\ln x \to +\infty$ (logaritmo in base $e > 1$).
  \item $\displaystyle\lim_{x\to 0^+}\left(x + \ln x\right) = -\infty$, perché $x \to 0$ e $\ln x \to -\infty$.
\end{itemize}
\end{esempio}

\section{Limiti delle funzioni trigonometriche}

\[
\lim_{x\to\pm\infty}\sin x, \qquad \lim_{x\to\pm\infty}\cos x, \qquad \lim_{x\to\pm\infty}\tan x \qquad\text{non esistono.}
\]
\[
\lim_{x\to\left(\frac{\pi}{2} + k\pi\right)^{\pm}} \tan x = \mp\infty, \quad k \in \mathbb{Z}, \qquad\qquad
\lim_{x\to\pm\infty}\arctan x = \pm\frac{\pi}{2}.
\]

\begin{proposizione}
Tutte le funzioni periodiche \emph{non} hanno limite per $x \to +\infty$ o $x \to -\infty$, tranne $f(x) = a$ con $a \in \mathbb{R}$, funzione costante, per cui $\displaystyle\lim_{x\to\pm\infty} f(x) = \lim_{x\to\pm\infty} a = a$.
\end{proposizione}

\begin{esempio}[Limite della tangente in $\frac{\pi}{2}$]
Sulla circonferenza goniometrica sia $A = (1, 0)$ e $Q$ l'intersezione tra la retta tangente in $A$ e la retta per l'origine che forma l'angolo $x$: $\tan x = \overline{AQ}$.
\begin{itemize}
  \item Se $0 < x < \frac{\pi}{2}$ e $x \to \frac{\pi}{2}$, il punto $Q$ sale lungo la tangente: $\tan x \to +\infty$.
  \item Se $\frac{\pi}{2} < x < \pi$ e $x \to \frac{\pi}{2}$, il punto $Q'$ scende, $\tan x' = \overline{AQ'}$ (con segno negativo): $\tan x \to -\infty$.
\end{itemize}
Quindi $\displaystyle\lim_{x\to\frac{\pi}{2}^-}\tan x = +\infty$ e $\displaystyle\lim_{x\to\frac{\pi}{2}^+}\tan x = -\infty$.
\end{esempio}

\end{document}
